\documentclass[conference]{IEEEtran}
\IEEEoverridecommandlockouts

\usepackage{cite}
\usepackage{amsmath,amssymb}
\usepackage{graphicx}
\usepackage{xcolor}
\usepackage{url}
% Local: figures live in runs/workshop3/figures/  Overleaf: figures/ (see scripts/pack_report.py)
\graphicspath{{../runs/workshop3/figures/}{figures/}}

% ---------------------------------------------------------------------------
% Helper macros
% \todo{...}  -> red text for things the authors must fill in
% \figph[width]{description}{file} -> inserts figures/<file>.png when it exists
%               (generate them with: python src/w3/make_report_assets_w3.py);
%               otherwise it shows a framed placeholder so the document compiles.
% generated/macros.tex -> every number of the text, written by the same script.
% ---------------------------------------------------------------------------
\newcommand{\todo}[1]{\textcolor{red}{\textbf{[TODO: #1]}}}
\newcommand{\figph}[3][\columnwidth]{%
  \includegraphics[width=#1]{#3.png}}
% Generated by src/w3/make_report_assets_w3.py -- do not edit by hand
\newcommand{\RandMeanReturn}{23.7}
\newcommand{\RandStdReturn}{12.2}
\newcommand{\RandShapedReturn}{23.7}
\newcommand{\RandMinLen}{9}
\newcommand{\RandMedianLen}{19.0}
\newcommand{\RandMaxLen}{63}
\newcommand{\RandLeftPct}{48.0}
\newcommand{\RandRightPct}{52.0}
\newcommand{\RandThetaFailPct}{100}
\newcommand{\ParityDiff}{0}
\newcommand{\BestConfig}{reference}
\newcommand{\BestSeed}{99}
\newcommand{\BestEpisode}{280}
\newcommand{\BestValReturn}{500.0}
\newcommand{\BestRunSteps}{36{,}500}
\newcommand{\NParams}{17{,}410}
\newcommand{\BestCfgBestVal}{500.0}
\newcommand{\BestCfgFinalVal}{187.9}
\newcommand{\TrainEpisodes}{300}
\newcommand{\MeanRunTime}{69}
\newcommand{\MinFinalVal}{96.2}
\newcommand{\MaxFinalVal}{207.1}
\newcommand{\StepsHalfMean}{6{,}246}
\newcommand{\StepsSolvedMean}{9{,}178}
\newcommand{\RunsSolved}{1/3}
\newcommand{\WOneDatasetSteps}{35{,}874}
\newcommand{\WOneTrainRows}{24{,}782}
\newcommand{\BestCheckpointSteps}{15{,}674}
\newcommand{\TestRandMean}{20.7}
\newcommand{\TestRandStd}{9.2}
\newcommand{\TestRandSuccess}{0}
\newcommand{\TestRandDanger}{23.67}
\newcommand{\TestRandDangerWTwo}{24.59}
\newcommand{\TestRandTheta}{4.01}
\newcommand{\TestRandAbsX}{0.05}
\newcommand{\TestRandXFail}{0}
\newcommand{\NoisyRandMean}{21.0}
\newcommand{\NoisyRandSuccess}{0}
\newcommand{\NoisyRandXFail}{0}
\newcommand{\NoisyRandThetaFail}{100}
\newcommand{\NoisyRandDanger}{22.60}
\newcommand{\NoisyRandDangerWTwo}{23.18}
\newcommand{\TestAltMean}{36.0}
\newcommand{\TestAltStd}{13.5}
\newcommand{\TestAltSuccess}{0}
\newcommand{\TestAltDanger}{14.42}
\newcommand{\TestAltDangerWTwo}{14.53}
\newcommand{\TestAltTheta}{3.95}
\newcommand{\TestAltAbsX}{0.05}
\newcommand{\TestAltXFail}{0}
\newcommand{\NoisyAltMean}{34.1}
\newcommand{\NoisyAltSuccess}{0}
\newcommand{\NoisyAltXFail}{0}
\newcommand{\NoisyAltThetaFail}{100}
\newcommand{\NoisyAltDanger}{14.53}
\newcommand{\NoisyAltDangerWTwo}{15.30}
\newcommand{\TestAngleMean}{42.1}
\newcommand{\TestAngleStd}{9.8}
\newcommand{\TestAngleSuccess}{0}
\newcommand{\TestAngleDanger}{18.79}
\newcommand{\TestAngleDangerWTwo}{10.17}
\newcommand{\TestAngleTheta}{4.74}
\newcommand{\TestAngleAbsX}{0.07}
\newcommand{\TestAngleXFail}{0}
\newcommand{\NoisyAngleMean}{40.8}
\newcommand{\NoisyAngleSuccess}{0}
\newcommand{\NoisyAngleXFail}{0}
\newcommand{\NoisyAngleThetaFail}{100}
\newcommand{\NoisyAngleDanger}{19.06}
\newcommand{\NoisyAngleDangerWTwo}{10.34}
\newcommand{\TestPDMean}{500.0}
\newcommand{\TestPDStd}{0.0}
\newcommand{\TestPDSuccess}{100}
\newcommand{\TestPDDanger}{0.00}
\newcommand{\TestPDDangerWTwo}{0.00}
\newcommand{\TestPDTheta}{0.22}
\newcommand{\TestPDAbsX}{0.37}
\newcommand{\TestPDXFail}{0}
\newcommand{\NoisyPDMean}{464.0}
\newcommand{\NoisyPDSuccess}{64}
\newcommand{\NoisyPDXFail}{36}
\newcommand{\NoisyPDThetaFail}{0}
\newcommand{\NoisyPDDanger}{0.42}
\newcommand{\NoisyPDDangerWTwo}{0.17}
\newcommand{\TestWOneMean}{492.7}
\newcommand{\TestWOneStd}{31.0}
\newcommand{\TestWOneSuccess}{91}
\newcommand{\TestWOneDanger}{0.75}
\newcommand{\TestWOneDangerWTwo}{0.28}
\newcommand{\TestWOneTheta}{3.46}
\newcommand{\TestWOneAbsX}{0.47}
\newcommand{\TestWOneXFail}{6}
\newcommand{\NoisyWOneMean}{438.5}
\newcommand{\NoisyWOneSuccess}{66}
\newcommand{\NoisyWOneXFail}{7}
\newcommand{\NoisyWOneThetaFail}{27}
\newcommand{\NoisyWOneDanger}{1.48}
\newcommand{\NoisyWOneDangerWTwo}{0.47}
\newcommand{\TestWTwoMean}{500.0}
\newcommand{\TestWTwoStd}{0.0}
\newcommand{\TestWTwoSuccess}{100}
\newcommand{\TestWTwoDanger}{0.00}
\newcommand{\TestWTwoDangerWTwo}{0.00}
\newcommand{\TestWTwoTheta}{1.92}
\newcommand{\TestWTwoAbsX}{0.52}
\newcommand{\TestWTwoXFail}{0}
\newcommand{\NoisyWTwoMean}{499.2}
\newcommand{\NoisyWTwoSuccess}{99}
\newcommand{\NoisyWTwoXFail}{1}
\newcommand{\NoisyWTwoThetaFail}{0}
\newcommand{\NoisyWTwoDanger}{0.01}
\newcommand{\NoisyWTwoDangerWTwo}{0.00}
\newcommand{\TestDQNMean}{498.1}
\newcommand{\TestDQNStd}{11.7}
\newcommand{\TestDQNSuccess}{97}
\newcommand{\TestDQNDanger}{9.64}
\newcommand{\TestDQNDangerWTwo}{9.54}
\newcommand{\TestDQNTheta}{6.07}
\newcommand{\TestDQNAbsX}{0.42}
\newcommand{\TestDQNXFail}{3}
\newcommand{\NoisyDQNMean}{370.9}
\newcommand{\NoisyDQNSuccess}{49}
\newcommand{\NoisyDQNXFail}{4}
\newcommand{\NoisyDQNThetaFail}{47}
\newcommand{\NoisyDQNDanger}{10.09}
\newcommand{\NoisyDQNDangerWTwo}{9.44}
\newcommand{\SeedsTestMean}{487.6}
\newcommand{\SeedsTestStd}{16.2}
\newcommand{\SeedsNoisyMean}{450.3}
\newcommand{\SeedsNoisyStd}{56.7}
\newcommand{\SeedsTestSuccess}{95}
\newcommand{\SeedsNoisySuccess}{80}
\newcommand{\NoiseStd}{2}
\newcommand{\WOneThreshold}{0.869}
\newcommand{\WTwoThreshold}{0.892}
\newcommand{\LR}{0.001}
\newcommand{\DiscountGamma}{0.99}
\newcommand{\EpsDecay}{0.99}
\newcommand{\EpsMin}{0.01}
\newcommand{\Buffer}{50{,}000}
\newcommand{\Batch}{64}
\newcommand{\LearnStarts}{1{,}000}
\newcommand{\TargetUpdate}{500}
\newcommand{\GradClip}{10}
\newcommand{\ShapingK}{1}
\newcommand{\EvalEvery}{10}
\newcommand{\WOneFOne}{0.860}
\newcommand{\WOneRecall}{0.886}
\newcommand{\WOnePrecision}{0.834}
\newcommand{\WOnePRAUC}{0.935}
\newcommand{\WTwoFOne}{0.865}
\newcommand{\WTwoRecall}{0.924}
\newcommand{\WTwoPrecision}{0.812}
\newcommand{\WTwoPRAUC}{0.956}
\newcommand{\WTwoWindow}{16}


\begin{document}

\title{Reinforcement Learning for Failure Avoidance in CartPole: a DQN Agent Grounded in Supervised Failure Prediction}

\author{\IEEEauthorblockN{Miguel Angel Torres Ramirez}
\IEEEauthorblockA{Dept. of Computer Engineering\\
Universidad Distrital Francisco Jos\'e de Caldas\\
Email: migatorresr@udistrital.edu.co}
\and
\IEEEauthorblockN{Cristian David Romero Gil}
\IEEEauthorblockA{Dept. of Computer Engineering\\
Universidad Distrital Francisco Jos\'e de Caldas\\
Email: cdromerog@udistrital.edu.co}
\and
\IEEEauthorblockN{Omar Esteban Pe\~na Sanchez}
\IEEEauthorblockA{Dept. of Computer Engineering\\
Universidad Distrital Francisco Jos\'e de Caldas\\
Email: oepenas@udistrital.edu.co}}

\maketitle

\begin{abstract}
In previous workshops we predicted, from labelled CartPole trajectories, whether an
episode was about to fail. This report turns the same system into a control problem:
an agent pushes the cart left or right at every step and must keep the pole upright
for as long as possible. We formulate the task as a Markov Decision Process, implement
a custom Gymnasium environment whose dynamics match CartPole-v1 exactly, add a
potential-based reward shaping built from the Workshop~1 finding that the distance to
the failure limits predicts failures, and train a Deep Q-Network implemented from
scratch. Six configurations, each changing one hyperparameter, were trained with three
seeds. The selected agent survives \TestDQNMean\ of 500 steps on average on 100
held-out initial states (\TestDQNSuccess\% success), against \TestRandMean\ for a random
policy, but it does not beat a hand-designed PD rule, and it operates closer to the
failure limits: the Workshop~1 and Workshop~2 models flag \TestDQNDanger\% and
\TestDQNDangerWTwo\% of its steps as dangerous, and its return falls to \NoisyDQNMean\
under actuator noise. The most robust controller was the Workshop~2 LSTM used as a
one-step look-ahead policy (\NoisyWTwoMean\ steps under noise), a supervised model
that, unlike DQN, relies on expert-generated labels and on a model of the dynamics.
Training curves show fast learning followed by policy degradation in every
configuration, which we diagnose and mitigate with validation-based checkpoint selection.
\end{abstract}

\begin{IEEEkeywords}
Reinforcement learning, Deep Q-Network, Markov decision process, reward shaping, CartPole, Gymnasium.
\end{IEEEkeywords}

%% ======================================================================
\section{Introduction}
Workshops~1 and~2 treated CartPole failures as a supervised problem: given the state
of the cart and the pole, decide whether the episode will fail within five steps. A
classifier of that kind tells \emph{when} a failure is coming, but not \emph{what to do}
about it. Choosing actions is a different problem: every push changes the states that
can be reached later, so the quality of a decision is only known after its consequences
unfold. This is the setting of reinforcement learning (RL), where an agent learns a
policy by interacting with an environment and receiving rewards \cite{sutton2018}.

RL fits our domain naturally for three reasons. First, the task is a sequence of
dependent decisions. Second, success is measured by a reward (the time the pole stays
up) rather than by labels: our dataset contains the actions of a noisy expert policy,
not the \emph{correct} action for each state. Third, the system can be simulated, so
interactions are cheap.

The contributions of this report are: (i) a complete MDP specification of the
control task, grounded in the findings of Workshop~1; (ii) a Gymnasium environment
validated against CartPole-v1 and with a random policy; (iii) a DQN agent
\cite{mnih2015} implemented from scratch, with a six-configuration, three-seed
hyperparameter search; and (iv) a test-set comparison against six baselines that need
no RL training, two of which turn the Workshop~1 classifier and the Workshop~2 LSTM into
policies, together with a sample-efficiency analysis.

Section~II defines the MDP, Section~III the environment, Section~IV the agent,
Section~V the hyperparameter search, Section~VI the test evaluation, Section~VII the
sample efficiency, Section~VIII discusses the results and Section~IX concludes.

%% ======================================================================
\section{MDP Formulation}

\subsection{Problem Statement}
The agent controls a cart that moves on a 4.8~m track and carries a pole hinged on
top. At each time step (0.02~s) it applies a horizontal force of $\pm 10$~N. The goal
is to keep the pole upright and the cart on the track for as long as possible, up to
500 steps (10~s). Table~\ref{tab:mdp} summarises the five components of the MDP.

\begin{table}[t]
\caption{MDP SPECIFICATION}
\label{tab:mdp}
\centering
\small
\begin{tabular}{|p{1.5cm}|p{5.9cm}|}
\hline
\textbf{Component} & \textbf{Definition} \\ \hline
State $s_t$ & $(x, \dot x, \theta, \dot\theta) \in \mathbb{R}^4$: cart position, cart velocity, pole angle, pole angular velocity \\ \hline
Actions $a_t$ & $\{0, 1\}$: push left ($-10$~N) or push right ($+10$~N) \\ \hline
Transition & cart-pole equations of motion, explicit Euler step of 0.02~s; deterministic given $(s_t, a_t)$; random initial state $s_0 \sim U(-0.05, 0.05)^4$ \\ \hline
Reward & $r_t = 1$ per step survived, plus optional shaping $F = \gamma\Phi(s_{t+1}) - \Phi(s_t)$ \\ \hline
Termination & failure if $|\theta| > 12^\circ$ or $|x| > 2.4$~m (terminal); truncation at 500 steps (not terminal) \\ \hline
\end{tabular}
\end{table}

\subsection{State Space}
The observation is the 4-dimensional physical state. It is Markovian: position,
velocity, angle and angular velocity completely determine the future evolution of the
system for a given sequence of forces. We did not feed the 16 engineered features of
Workshop~1 to the agent: they are functions of the current and past states, so they add
no information for a Markovian controller, and the Q-network can learn its own
features. The inputs are divided by fixed physical scales (the failure limits for $x$
and $\theta$, and typical magnitudes for the velocities); these constants come from the
environment definition, not from data.

\subsection{Action Set}
The action space is discrete with two actions: 0 = push left and 1 = push right. There
is no ``do nothing'' action, as in the original problem \cite{barto1983}: balancing
results from alternating pushes.

\subsection{Transition Dynamics}
The next state $s_{t+1} = f(s_t, a_t)$ follows the frictionless cart-pole equations
\cite{barto1983}. With $F = \pm 10$~N, $m_p = 0.1$~kg, $M = m_c + m_p = 1.1$~kg,
half-length $l = 0.5$~m and $g = 9.8$~m/s$^2$:
\begin{align}
\ddot\theta &= \frac{g\sin\theta - \cos\theta\,\big(F + m_p l \dot\theta^2 \sin\theta\big)/M}
                    {l\left(\frac{4}{3} - \frac{m_p\cos^2\theta}{M}\right)}, \\
\ddot x &= \frac{F + m_p l\left(\dot\theta^2\sin\theta - \ddot\theta\cos\theta\right)}{M},
\end{align}
followed by an explicit Euler update of the four variables with $\tau = 0.02$~s. The
transition is deterministic given the action. Stochasticity comes from the random
initial state and, in the robustness test of Section~VI, from Gaussian noise added to
the force.

\subsection{Reward Function}
The \emph{task reward} is $r_t = +1$ for every step the system survives, so the
undiscounted return equals the number of steps survived (maximum 500). It is dense in
time but says nothing about \emph{how close} the system is to failing: the only negative
signal is the loss of future rewards when the episode ends.

Workshop~1 showed that the distance to the failure limits is the best predictor of a
failure: adding it raised the ROC-AUC of the linear classifier from 0.542 to 0.981. We
encode that finding as a potential function
\begin{equation}
\Phi(s) = -\frac{k}{2}\left[\left(\frac{x}{2.4}\right)^2 + \left(\frac{\theta}{0.2095}\right)^2\right],
\end{equation}
which is 0 at the centred upright position and decreases towards the limits, and train
with the shaped reward $r_t + F_t$, where $F_t = \gamma\Phi(s_{t+1}) - \Phi(s_t)$ and
$\Phi(\text{terminal}) = 0$. Potential-based shaping provably leaves the optimal policy
unchanged \cite{ng1999}: over an episode the discounted shaping terms telescope to
approximately $-\Phi(s_0)$, a constant close to zero, so they can only change \emph{how fast} the
agent learns, not \emph{what} it learns. We use $k = \ShapingK$; the configuration
\texttt{no\_shaping} ($k = 0$) measures its effect. All agents are evaluated on the task
reward only. The relevant time horizon is long: with $\gamma = \DiscountGamma$ the
effective horizon $1/(1-\gamma)$ is about 100 steps, comparable to the time it takes a
drifting cart to leave the track.

\subsection{Episode Termination}
An episode \emph{terminates} (failure) when $|\theta| > 12^\circ$ or $|x| > 2.4$~m, the
same limits that defined the danger label of Workshop~1. It is \emph{truncated} after
500 steps: that is a time limit, not a failure, so the agent still bootstraps from the
last state; only terminations set the value of the next state to zero. Reaching the
time limit is a success.

%% ======================================================================
\section{Environment Implementation and Validation}

\subsection{Implementation}
\texttt{CartPoleControlEnv} (\texttt{src/w3/env.py}) subclasses
\texttt{gymnasium.Env} \cite{towers2024}. It declares
\texttt{action\_space = Discrete(2)} and a 4-dimensional \texttt{Box} observation
space, implements \texttt{reset()} (random initial state from the seeded generator)
and \texttt{step()} with the transition, the reward (task reward plus shaping) and the
termination and truncation logic. The \texttt{info} dictionary reports the task reward
and the outcome (\emph{failure by angle}, \emph{failure by position} or \emph{time
limit}). \texttt{render()} returns a text line or an RGB image drawn with NumPy
(Fig.~\ref{fig:render}).

\begin{figure}[t]
\centering
\figph{frames produced by env.render()}{fig-render}
\caption{Frames produced by \texttt{render()} along a test episode of the selected agent (red marks: failure limits $|x| = 2.4$~m).}
\label{fig:render}
\end{figure}

\subsection{Validation}
Three checks were run before any training. (1) Gymnasium's \texttt{check\_env}
confirmed the API (spaces, signatures and data types). (2) Our environment and
Gymnasium's CartPole-v1, reset with the same seeds and fed the same random actions for
20 episodes, produced a maximum state difference of \ParityDiff: the dynamics are
identical to those of the system that generated our dataset. (3) We ran 100 episodes
with a uniformly random policy (Fig.~\ref{fig:random}): the mean return was
\RandMeanReturn~$\pm$~\RandStdReturn\ (length between \RandMinLen\ and \RandMaxLen\
steps, median \RandMedianLen), both actions occurred (\RandLeftPct\% left,
\RandRightPct\% right) and \RandThetaFailPct\% of the episodes ended because the pole
fell. With $k = 1$ the mean shaped return was \RandShapedReturn, almost equal to the
task return, as expected from the telescoping property.

\begin{figure*}[t]
\centering
\figph[0.8\textwidth]{random-policy diagnostics}{fig-random-policy}
\caption{Random-policy diagnostics on 100 episodes: (a) episode-length distribution, (b) frequency of each action.}
\label{fig:random}
\end{figure*}

%% ======================================================================
\section{Agent: Deep Q-Network}

\subsection{Algorithm and Architecture}
DQN \cite{mnih2015} suits a discrete action space with Markovian dynamics. It
approximates the action-value function $Q(s, a; w)$, the expected discounted return of
taking action $a$ in state $s$ and acting greedily afterwards, and is trained toward the
Q-learning target \cite{watkins1992}
\begin{equation}
y = r + \gamma\,(1 - d)\max_{a'} Q(s', a'; w^-),
\end{equation}
where $d = 1$ only for a failure and $w^-$ are the weights of a target network. The
Q-network (implemented in PyTorch \cite{paszke2019}) is a multilayer perceptron
$4 \to 128 \to 128 \to 2$ with ReLU activations and a linear output, \NParams\
trainable parameters.

\subsection{Training Loop}
Each environment step stores the transition $(s, a, r, s', d)$ in a replay buffer of
\Buffer\ transitions; after \LearnStarts\ steps, every step samples a mini-batch of
\Batch\ transitions uniformly, which breaks the correlation between consecutive states.
The loss is the Huber loss between $Q(s, a; w)$ and $y$, minimised with Adam; the
gradient norm is clipped at \GradClip\ and the target network is a hard copy of the
online network every \TargetUpdate\ gradient steps. Exploration is
$\epsilon$-greedy: $\epsilon$ starts at 1 and is multiplied by the decay factor after
every episode, down to \EpsMin. Each run lasts \TrainEpisodes\ episodes; per-episode
return, length, $\epsilon$ and loss are logged to TensorBoard and CSV files in
\texttt{runs/workshop3/}. Table~\ref{tab:hparams} lists the reference values.

\begin{table}[t]
\caption{REFERENCE HYPERPARAMETERS}
\label{tab:hparams}
\centering
\small
\begin{tabular}{|l|l|}
\hline
\textbf{Hyperparameter} & \textbf{Value} \\ \hline
Learning rate (Adam) & \LR \\ \hline
Discount factor $\gamma$ & \DiscountGamma \\ \hline
$\epsilon$ decay per episode / minimum & \EpsDecay\ / \EpsMin \\ \hline
Weight decay (L2) & 0 \\ \hline
Shaping strength $k$ & \ShapingK \\ \hline
Replay buffer / batch size & \Buffer\ / \Batch \\ \hline
Learning starts & \LearnStarts\ steps \\ \hline
Target-network update & every \TargetUpdate\ updates \\ \hline
Gradient clipping (L2 norm) & \GradClip \\ \hline
Episodes / validation interval & \TrainEpisodes\ / \EvalEvery\ episodes \\ \hline
\end{tabular}
\end{table}

\subsection{Validation Loop and Checkpointing}
Initial states are split into three disjoint sets: training episodes use seeds above
$10^5$, validation uses 10 fixed seeds and the test uses 100 other fixed seeds.
Every \EvalEvery\ episodes the agent is evaluated greedily ($\epsilon = 0$) on the
validation episodes. We report the mean validation return and its moving average over
the last three evaluations (smoothed validation return), and keep the checkpoint with
the highest validation return of each run. The best run of the selected configuration
is saved as \texttt{checkpoints/workshop3/agent\_best.pt}.

%% ======================================================================
\section{Hyperparameter Search}

\subsection{Protocol}
Six configurations were trained with seeds 42, 7 and 99 (18 runs). Each one changes a
single factor with respect to \texttt{reference}, so that differences can be attributed
to that factor: a lower learning rate, a lower discount factor, slower exploration
decay, L2 weight decay, and no reward shaping. A configuration is selected with
validation results only: the highest mean (over seeds) best validation return, ties
broken by the final validation return (mean of the last five evaluations), which
measures stability at the end of training.

\begin{table*}[t]
\caption{HYPERPARAMETER SEARCH: VALIDATION RESULTS (MEAN $\pm$ STD OVER 3 SEEDS)}
\label{tab:sweep}
\centering
\small
\setlength{\tabcolsep}{4pt}
\begin{tabular}{|l|c|c|c|c|c|c|c|c|c|}
\hline
\textbf{Configuration} & \textbf{LR} & $\boldsymbol{\gamma}$ & \textbf{$\epsilon$ decay} & \textbf{L2} & $\boldsymbol{k}$ &
\textbf{Best val. return} & \textbf{Final val. return} & \textbf{Runs $\geq$ 475} & \textbf{Steps to 475} \\ \hline
\textbf{reference} & 0.001 & 0.99 & 0.99 & 0 & 1 & $500.0 \pm 0.0$ & $187.9 \pm 98.3$ & 1/3 & 9{,}178 \\ \hline
lr\_1e-4 & 0.0001 & 0.99 & 0.99 & 0 & 1 & $389.9 \pm 83.3$ & $129.1 \pm 55.4$ & 0/3 & -- \\ \hline
gamma\_0.95 & 0.001 & 0.95 & 0.99 & 0 & 1 & $500.0 \pm 0.0$ & $96.2 \pm 3.0$ & 2/3 & 16{,}902 \\ \hline
eps\_decay\_0.995 & 0.001 & 0.99 & 0.995 & 0 & 1 & $464.4 \pm 50.3$ & $204.2 \pm 72.1$ & 0/3 & -- \\ \hline
weight\_decay\_1e-4 & 0.001 & 0.99 & 0.99 & 0.0001 & 1 & $496.0 \pm 5.7$ & $207.1 \pm 62.4$ & 0/3 & -- \\ \hline
no\_shaping & 0.001 & 0.99 & 0.99 & 0 & 0 & $498.8 \pm 1.6$ & $192.4 \pm 20.2$ & 0/3 & -- \\ \hline

\end{tabular}
\end{table*}

\begin{figure*}[t]
\centering
\figph[0.95\textwidth]{validation return vs environment steps for every configuration and seed}{fig-training-curves}
\caption{Validation return (10 episodes, greedy policy) as a function of environment steps for every configuration and seed. Dashed line: 475, the CartPole-v1 solving threshold.}
\label{fig:curves}
\end{figure*}

\subsection{Results}
Table~\ref{tab:sweep} and Fig.~\ref{fig:curves} summarise the 18 runs. ``Runs $\geq$
475'' counts the runs whose \emph{smoothed} validation return reached 475, which
requires three consecutive good evaluations. Every configuration shows measurable
learning: the validation return rises from about 10 to several hundred within the
first few thousand steps. The selected configuration is \texttt{\BestConfig}, whose
three seeds all reached a validation return of \BestCfgBestVal; its best run (seed
\BestSeed, episode \BestEpisode) is \texttt{agent\_best.pt}.

The most striking pattern is shared by all configurations: the final validation return
(\MinFinalVal\ to \MaxFinalVal) is far below the best one. The agents learn a good policy
and then lose it, a known instability of DQN caused by bootstrapping from a moving
target with function approximation. Fig.~\ref{fig:explore}(b) shows how it happens for
the reference configuration: during the first two thousand steps the episodes end
because the pole falls; once the agent has learned to hold the pole, almost every
validation episode that does not reach 500 steps ends because the cart leaves the
track, including those of the degraded policies late in training (beyond about 22
thousand steps the curve comes from a single seed, the only run that lasted that long).
The degraded policy keeps the pole up but lets the cart drift, a failure
whose cost appears about 100 steps later, at the edge of the effective horizon of
$\gamma = 0.99$. Selecting checkpoints on validation is what makes the final agent
usable despite this degradation.

Regarding individual factors: with the lower learning rate none of the seeds reached a
validation return of 500, so its best return is the lowest of the search; $\gamma = 0.95$ reached 500 in all seeds
but degraded the most, to a final return of about 96, consistent with a horizon of about
20 steps that cannot see the drift failure; slower exploration decay and weight decay
did not change the picture; and removing the shaping produced practically the same best
and final returns as the reference. With $k = 1$ the shaping had no measurable effect,
which is compatible with its theoretical role (it cannot change the optimal policy) but
means it did not speed up learning either.

\begin{figure*}[t]
\centering
\figph[0.95\textwidth]{epsilon schedules and how validation episodes end}{fig-exploration}
\caption{(a) Exploration rate per episode for the two decay factors. (b) Reference configuration: share of validation episodes that end by cart position or by pole angle during training (mean of three seeds).}
\label{fig:explore}
\end{figure*}

%% ======================================================================
\section{Test Evaluation and Comparison}

\subsection{Protocol and Baselines}
The test set contains 100 initial states never used for training or selection. It is
evaluated in two conditions: \emph{nominal} (the training dynamics) and \emph{noisy},
with Gaussian actuator noise of standard deviation \NoiseStd~N added to every push, a
perturbation the agent never saw. The DQN agent is compared with six policies that
need no RL training:
\begin{itemize}
\item \textbf{Random:} uniform random actions.
\item \textbf{Alternating (fixed):} left, right, left\dots, a fixed open-loop strategy.
\item \textbf{Angle rule (greedy):} push towards the side the pole leans.
\item \textbf{PD rule (heuristic):} push towards the sign of $\theta + 0.5\,\dot\theta$; the gain was chosen on the validation episodes.
\item \textbf{W1 classifier (look-ahead):} the Workshop~1 Gradient Boosting classifier used as a policy. For each action, the known dynamics give the next state and the classifier scores its probability of danger; the action with the lower probability is chosen. Unlike DQN, it uses a model of the dynamics.
\item \textbf{W2 LSTM (look-ahead):} the same rule with the Workshop~2 model, a 2-layer LSTM that reads the last \WTwoWindow\ steps of the episode (states and the actions actually taken). The candidate step carries the action that leads to it, because the action that will be taken there is not known yet. The model is loaded from \texttt{checkpoints/workshop2/model\_best.pt} without retraining; on the Workshop~2 test episodes our bridge reproduces its reported test F1 exactly.
\end{itemize}

Besides the return, we report the success rate (500 steps reached), the share of
failures by cart position, the mean $|\theta|$, and two \emph{danger rates}: the share
of steps that the Workshop~1 classifier (threshold \WOneThreshold) and the Workshop~2
LSTM (threshold \WTwoThreshold), each with the threshold chosen on its own validation
set, flag as being within five steps of a failure. These metrics link the three
workshops: they measure how much time each policy spends in states that the supervised
models consider dangerous.

\begin{table*}[t]
\caption{TEST RESULTS ON 100 HELD-OUT EPISODES (RETURN = STEPS SURVIVED, MAXIMUM 500)}
\label{tab:test}
\centering
\small
\setlength{\tabcolsep}{3.5pt}
\begin{tabular}{|l|c|c|c|c|c|c||c|c|}
\hline
& \multicolumn{6}{c||}{\textbf{Nominal test}} & \multicolumn{2}{c|}{\textbf{Noisy test (\NoiseStd~N)}} \\ \hline
\textbf{Policy} & \textbf{Return} & \textbf{Succ. (\%)} & \textbf{Fail $x$ (\%)} & \textbf{W1 dang. (\%)} & \textbf{W2 dang. (\%)} & \textbf{$|\theta|$ (deg)} & \textbf{Return} & \textbf{Succ. (\%)} \\ \hline
Random & $20.7 \pm 9.2$ & 0 & 0 & 23.67 & 24.59 & 4.01 & $21.0 \pm 11.5$ & 0 \\ \hline
Alternating (fixed) & $36.0 \pm 13.5$ & 0 & 0 & 14.42 & 14.53 & 3.95 & $34.1 \pm 14.2$ & 0 \\ \hline
Angle rule (greedy) & $42.1 \pm 9.8$ & 0 & 0 & 18.79 & 10.17 & 4.74 & $40.8 \pm 10.7$ & 0 \\ \hline
PD rule (heuristic) & $500.0 \pm 0.0$ & 100 & 0 & 0.00 & 0.00 & 0.22 & $464.0 \pm 65.0$ & 64 \\ \hline
W1 classifier (look-ahead) & $492.7 \pm 31.0$ & 91 & 6 & 0.75 & 0.28 & 3.46 & $438.5 \pm 107.8$ & 66 \\ \hline
W2 LSTM (look-ahead) & $500.0 \pm 0.0$ & 100 & 0 & 0.00 & 0.00 & 1.92 & $499.2 \pm 8.2$ & 99 \\ \hline
\textbf{DQN (agent\_best.pt)} & $498.1 \pm 11.7$ & 97 & 3 & 9.64 & 9.54 & 6.07 & $370.9 \pm 158.7$ & 49 \\ \hline
DQN, 3 seeds of reference & $487.6 \pm 16.2$ & 95 & -- & -- & -- & -- & $450.3 \pm 56.7$ & 80 \\ \hline

\end{tabular}
\end{table*}

\begin{figure}[t]
\centering
\figph{test return of every policy}{fig-test-comparison}
\caption{Mean test return ($\pm$ one standard deviation) of every policy in the nominal and noisy test conditions.}
\label{fig:test}
\end{figure}

\subsection{Results}
Table~\ref{tab:test} and Fig.~\ref{fig:test} show the results. In the nominal
condition the selected agent survives \TestDQNMean~$\pm$~\TestDQNStd\ steps with
\TestDQNSuccess\% success, far above the random (\TestRandMean), alternating
(\TestAltMean) and angle-rule (\TestAngleMean) policies. Across the three seeds of the
selected configuration the mean is \SeedsTestMean~$\pm$~\SeedsTestStd, so the result does
not depend on a lucky seed. The agent did not outperform the PD rule or the W2-guided
policy, which both solve every nominal test episode, and it is slightly above the W1
classifier used as a policy (\TestWOneMean, \TestWOneSuccess\% success).

The danger rate and the trajectories reveal \emph{how} the agent survives. The PD rule
keeps the pole almost vertical (mean $|\theta|$ of \TestPDTheta$^\circ$) and spends
\TestPDDanger\% of its steps in danger. The selected agent swings the pole between about
$-11^\circ$ and $+11^\circ$, close to the $12^\circ$ limit (Fig.~\ref{fig:traj}). The two
supervised models, trained independently and with different inputs, agree on this: the
Workshop~1 classifier flags \TestDQNDanger\% of its steps as dangerous and the Workshop~2
LSTM \TestDQNDangerWTwo\%. Both policies maximise the task reward equally, but the agent
has a much smaller safety margin, and the noisy test exposes it: its return drops to
\NoisyDQNMean\ with \NoisyDQNSuccess\% success, and \NoisyDQNThetaFail\% of the episodes
end because the pole falls. The PD rule (\NoisyPDMean) and the W1-guided policy
(\NoisyWOneMean) degrade less: the PD rule loses episodes through cart drift
(\NoisyPDXFail\% end at the track limit), the W1-guided policy mostly because the pole
falls (\NoisyWOneThetaFail\%). The W2-guided policy is the
most robust of all: \NoisyWTwoMean\ steps with \NoisyWTwoSuccess\% success under noise,
keeping the pole within a mean $|\theta|$ of \TestWTwoTheta$^\circ$ in the nominal test.

Table~\ref{tab:seeds} shows that this fragility belongs to the selected checkpoint,
not to the configuration: the other two seeds have both danger rates below 0.2\% and
keep most of their performance under noise. Among these three agents the danger rates,
measured in the nominal condition, identified the fragile one, whereas the validation
return, equal to 500 for the three, could not. We note this as an observation on three
agents, not as a general result.

\begin{table}[t]
\caption{SELECTED CONFIGURATION: TEST RESULTS PER TRAINING SEED}
\label{tab:seeds}
\centering
\small
\setlength{\tabcolsep}{3pt}
\begin{tabular}{|l|c|c|c|c|c||c|c|}
\hline
& \multicolumn{5}{c||}{\textbf{Nominal}} & \multicolumn{2}{c|}{\textbf{Noisy}} \\ \hline
\textbf{Seed} & \textbf{Ret.} & \textbf{Succ.} & \textbf{W1 d.} & \textbf{W2 d.} & $|\theta|$ & \textbf{Ret.} & \textbf{Succ.} \\ \hline
42 & 464.7 & 88 & 0.14 & 0.00 & 1.06 & 479.9 & 92 \\ \hline
7 & 500.0 & 100 & 0.02 & 0.00 & 5.26 & 500.0 & 100 \\ \hline
99 (agent\_best) & 498.1 & 97 & 9.64 & 9.54 & 6.07 & 370.9 & 49 \\ \hline

\end{tabular}
\end{table}

\subsection{Comparison with Workshops 1 and 2}
The supervised models of the previous workshops measure a different quantity
(classification of danger states), so their metrics cannot be placed in the same
column as a return; Table~\ref{tab:w12} relates both views. On the common test
episodes of Workshops~1 and~2, the Workshop~2 LSTM was the better predictor
(F1 \WTwoFOne\ versus \WOneFOne, recall \WTwoRecall\ versus \WOneRecall, PR-AUC
\WTwoPRAUC\ versus \WOnePRAUC), and as a look-ahead policy it was also the better
controller (\TestWTwoMean\ versus \TestWOneMean\ nominal, \NoisyWTwoMean\ versus
\NoisyWOneMean\ with noise). A plausible reason, which we did not test separately, is
that its 16-step window lets it judge where the system is heading, while the W1 features
look back only three steps. Both supervised controllers rely on two resources the DQN
agent did not have: labels produced by an expert policy and a model of the dynamics to
simulate each action.

\begin{table}[t]
\caption{SUPERVISED MODELS AS PREDICTORS (W1/W2 TEST SET) AND AS LOOK-AHEAD POLICIES (W3 TEST)}
\label{tab:w12}
\centering
\small
\setlength{\tabcolsep}{2.5pt}
\begin{tabular}{|l|c|c|c||c|c|}
\hline
& \multicolumn{3}{c||}{\textbf{Prediction}} & \multicolumn{2}{c|}{\textbf{Control (return)}} \\ \hline
\textbf{Model} & \textbf{F1} & \textbf{Recall} & \textbf{PR-AUC} & \textbf{Nominal} & \textbf{Noisy} \\ \hline
W1 Gradient Boosting & \WOneFOne & \WOneRecall & \WOnePRAUC & \TestWOneMean & \NoisyWOneMean \\ \hline
W2 LSTM & \WTwoFOne & \WTwoRecall & \WTwoPRAUC & \TestWTwoMean & \NoisyWTwoMean \\ \hline
DQN (this work) & -- & -- & -- & \TestDQNMean & \NoisyDQNMean \\ \hline
\end{tabular}
\end{table}

\subsection{Sample Trajectories}
Fig.~\ref{fig:traj} shows four test episodes of the selected agent: pole angle, cart
position, actions (push right in blue) and cumulative return. The agent alternates
pushes in bursts, producing a regular oscillation of the pole while the cart stays
within about one meter of the centre; the return grows by one per step until the
time limit.

\begin{figure*}[t]
\centering
\figph[0.95\textwidth]{four test episodes of the selected agent}{fig-trajectories}
\caption{Four test episodes of the selected agent: pole angle, cart position, action taken at each step and cumulative task return.}
\label{fig:traj}
\end{figure*}

%% ======================================================================
\section{Sample Efficiency}
Fig.~\ref{fig:sample} plots the smoothed validation return of the selected
configuration against the number of environment steps. On average its seeds first
reached a smoothed return of 195 after \StepsHalfMean\ steps; the solving level of 475
was reached by \RunsSolved\ seeds (after \StepsSolvedMean\ steps), and the best
checkpoints were saved after \BestCheckpointSteps\ steps on average.

\begin{figure*}[t]
\centering
\figph[0.95\textwidth]{learning curve of the selected configuration against the baselines}{fig-sample-efficiency}
\caption{Sample efficiency: smoothed validation return of the selected configuration (mean and min--max of three seeds) versus environment steps, with the test return of the baselines (PD rule and W2 LSTM share the line at 500) and the number of environment steps stored in the dataset of Workshops~1 and~2.}
\label{fig:sample}
\end{figure*}

For comparison, the dataset of Workshops~1 and~2 contains \WOneDatasetSteps\ environment
steps, of which \WOneTrainRows\ labelled rows were used to train both supervised models.
The RL agent therefore reached competent control with fewer environment interactions than
were needed to build the supervised dataset, although the supervised models, once turned
into controllers, ended up more robust. The comparison has two caveats. The supervised
models learn to \emph{predict} failures, and turning them into a controller required the
dynamics model used by the look-ahead policy; and the dataset itself was produced by an
expert policy, a source of knowledge that the RL agent did not have. Conversely, the PD
rule needs no data at all, only knowledge of the physics.

%% ======================================================================
\section{Discussion}

\emph{Did RL outperform supervised learning?} No. In the nominal environment the RL agent
matches the W1-guided policy and comes close to the PD rule and the W2-guided policy,
learning from rewards alone and without a model of the dynamics. Under actuator noise it
was the least robust of these four, and the W2 LSTM with look-ahead was the best
controller overall. The comparison is not like for like: the supervised controllers used
expert-generated labels and the simulator as a dynamics model, the DQN agent only its own
interactions. On a system this simple and well understood, a hand-designed rule also
remains a strong baseline.

\emph{Key challenges.} The main one was training instability: all configurations reached
a good policy and then degraded, so the final agent depends on validation-based
checkpoint selection (Fig.~\ref{fig:curves}). The second is that the task reward does
not reward safety margins: surviving with the pole near the limit earns the same
reward as keeping it vertical, which explains the oscillating behaviour of the selected
agent and its fragility under noise. Validating only on nominal episodes could not detect
this; adding noise to the validation episodes, or the W1/W2 danger rates, to the
selection criterion is a natural improvement.

\emph{Exploration versus exploitation.} With a decay of \EpsDecay\ per episode,
$\epsilon$ falls below 0.1 after about 230 episodes; the slower decay of 0.995 keeps it
above 0.2 for the whole run (Fig.~\ref{fig:explore}a). Slower decay produced shorter,
noisier training episodes, and therefore fewer environment steps per run, without
improving the best validation return. Late in training, when exploration is low, the policy drifted
toward the cart-position failure mode, which suggests that low exploration together with
bootstrapping lets errors in rarely visited states (near the track edges) go uncorrected.

\emph{Reward shaping.} The potential-based shaping with $k = 1$ did not change the
results measurably. Because it cannot change the optimal policy, it also cannot make the
agent prefer larger safety margins; a reward that explicitly penalises proximity to the
limits would change the task, and therefore the optimal policy, which is a design
decision rather than a tuning one.

\emph{Limitations.} Three seeds per configuration and ten validation episodes give noisy
estimates; hyperparameters were varied one at a time, so interactions between them
were not explored; and the comparison with the supervised models is indirect because
they solve a prediction task. Double DQN, soft target updates or prioritised replay
could reduce the degradation.

%% ======================================================================
\section{Conclusions}
We formulated CartPole control as an MDP grounded in our previous failure-prediction
work, implemented and validated a Gymnasium environment identical in dynamics to
CartPole-v1, and trained a from-scratch DQN agent with a six-configuration,
three-seed hyperparameter search. The agent learned measurably, surviving
\TestDQNMean\ of 500 steps on held-out episodes against \TestRandMean\ for a random
policy, but it did not beat a PD rule, and it operates with a small safety margin that
makes it fragile under actuator noise. The Workshop~1 and Workshop~2 models, reused as
danger monitors, revealed that margin, and the Workshop~2 LSTM, reused as a look-ahead
policy, was the most robust controller. Future work includes noise-aware validation,
reward functions that value safety margins explicitly, and stabilised variants of DQN.

\section*{Reproducibility}
Code, configurations, logs and checkpoints are available at
\url{https://github.com/Tonyaa11/Cartpole} (code in \path{src/w3/}, configurations in
\path{configs/sweep.json}, logs in \path{runs/workshop3/}, checkpoints in
\path{checkpoints/workshop3/}). All numbers and figures of this report are generated
from the logs by \path{src/w3/make_report_assets_w3.py}; seeds are fixed (training
42, 7, 99; validation 10000--10009; test 20000--20099). The Workshop~1 classifier is
refitted with scikit-learn 1.6.1, the version of the Workshop~1 report. Neural-network
training can differ slightly across library versions and hardware.

\begin{thebibliography}{00}
\bibitem{sutton2018} R. S. Sutton and A. G. Barto, \emph{Reinforcement Learning: An Introduction}, 2nd ed. Cambridge, MA, USA: MIT Press, 2018.
\bibitem{mnih2015} V. Mnih et al., ``Human-level control through deep reinforcement learning,'' \emph{Nature}, vol. 518, no. 7540, pp. 529--533, 2015.
\bibitem{watkins1992} C. J. C. H. Watkins and P. Dayan, ``Q-learning,'' \emph{Mach. Learn.}, vol. 8, no. 3--4, pp. 279--292, 1992.
\bibitem{barto1983} A. G. Barto, R. S. Sutton, and C. W. Anderson, ``Neuronlike adaptive elements that can solve difficult learning control problems,'' \emph{IEEE Trans. Syst., Man, Cybern.}, vol. SMC-13, no. 5, pp. 834--846, 1983.
\bibitem{ng1999} A. Y. Ng, D. Harada, and S. Russell, ``Policy invariance under reward transformations: Theory and application to reward shaping,'' in \emph{Proc. 16th Int. Conf. Mach. Learn. (ICML)}, 1999, pp. 278--287.
\bibitem{towers2024} M. Towers et al., ``Gymnasium: A standard interface for reinforcement learning environments,'' arXiv:2407.17032, 2024.
\bibitem{paszke2019} A. Paszke et al., ``PyTorch: An imperative style, high-performance deep learning library,'' in \emph{Adv. Neural Inf. Process. Syst. (NeurIPS)}, vol. 32, 2019.
\end{thebibliography}

\end{document}
